% RESUMO
Este é o template para confecção do artigo do Trabalho de Conclusão dos cursos de Ciência da Computação e Sistemas de Informação do Instituto de Matemática e Computação da Universidade Federal de Itajubá. Serão dadas dicas para confecção do artigo, além de exemplos de utilização de recursos do \LaTeX~ (ex. tabelas, figuras etc). Para mais informações leia ``A técnica da Escrita Científica'' de \cite{oliveira2015} disponível em \url{www.scielo.br/j/rbef/a/VsMNNbVGBzRkSwHQgKVF3cr/?lang=pt}.

\begin{center}
    \textcolor{red}{\LARGE\textbf{\underline{ATENÇÃO}}}
\end{center} 

Atente-se para a correta confecção deste artigo pois ele contém SEU trabalho e leva SEU nome. Mesmo que o orientador tenha sido o mentor do tema do trabalho, entenda que uma vez que se propôs a executá-lo você se torna o responsável pelo desenvolvimento e devida conclusão. Seja ético na confecção do seu artigo, referenciando corretamente os autores que contribuíram com conceitos e ideias para o seu trabalho. \\

A reprodução literal de um texto ou figura sem a correta atribuição da fonte por meio de citação e referência constitui crime de plágio conforme Lei 9.610 de 1998, que regula os direitos autorais artísticos e acadêmicos. Segundo o Art. 184 do Código Penal, o plagiador está sujeito a multa e até a detenção de 3 meses a 1 ano. Uma constatação de plágio mesmo após a formatura ou obtenção de título, pode acarretar na cassação do diploma e afetar negativamente seu currículo. A responsabilização pelo crime de plágio também está prevista na ``Norma Disciplinar do Corpo Discente da UNIFEI'' disponível em \url{https://drive.google.com/file/d/1BFfvx8JemOOi56GAxb03Kfi9jTM0342D/view}.